\section{Methodology}

\subsection{Evaluation protocol}

We use a \textbf{chronological 70/10/20 split} on rating timestamps:
\begin{itemize}
  \item Train: 700{,}146 rows (earliest 70\% of timeline)
  \item Validation: 100{,}020 rows (next 10\%)
  \item Test: 200{,}043 rows (final 20\%)
\end{itemize}

Hyperparameters are selected using validation RMSE; test metrics are reported once per model family.
This avoids leakage from future ratings into training features or tuning decisions.

\subsection{Feature engineering}

From raw tables we compute:
\begin{itemize}
  \item \texttt{user\_rating\_count}, \texttt{user\_avg\_rating}
  \item \texttt{movie\_rating\_count}, \texttt{movie\_avg\_rating}
  \item \texttt{release\_year} (parsed from movie title)
  \item User fields: \texttt{age}, \texttt{occupation}
\end{itemize}

Implementation: \texttt{src/features.py}, orchestrated in the capstone notebook.

\subsection{Machine learning models}

We train four regressors on the tabular feature matrix:
\begin{enumerate}
  \item \textbf{Baseline:} global mean rating on the training set.
  \item \textbf{Linear regression} (sklearn defaults).
  \item \textbf{Random forest} (\texttt{n\_estimators=100}, \texttt{max\_depth=10}).
  \item \textbf{Gradient boosting} (\texttt{n\_estimators=200}, \texttt{max\_depth=5},
        \texttt{learning\_rate=0.1}).
\end{enumerate}

ML hyperparameter search (\texttt{src/tune\_ml.py}) uses a fast-pass subsample of train/validation for
grid search, then refits the best configuration on the full training split.

\subsection{Neural collaborative filtering}

The NCF model (\texttt{src/dl\_model.py}) embeds users and movies, concatenates embeddings with dense
tabular features, and passes them through fully connected layers with dropout.
Baseline settings include embedding dimension 32, learning rate 0.001, and 10 training epochs.

\textbf{Optuna tuning} (\texttt{src/tune\_dl.py}) runs 50 trials with 3 epochs per trial on validation RMSE.
The winning configuration is retrained in \texttt{src/post\_analysis.py} to produce
\texttt{ncf\_tuned\_best.pt} and explainability artifacts.

\subsection{Metrics}

Primary offline metrics:
\begin{itemize}
  \item \textbf{RMSE} --- $\sqrt{\frac{1}{n}\sum_i (y_i - \hat{y}_i)^2}$; penalizes large errors.
  \item \textbf{MAE} --- mean absolute error; more robust to outliers.
\end{itemize}

Ranking metrics (Precision@K, Recall@K) are defined in \texttt{src/evaluation.py} for extension; the
capstone notebook focuses on rating regression consistent with the training objective.

\subsection{Reproducibility}

All steps are runnable from \texttt{notebooks/MovieMind\_capstone.ipynb}.
Reviewers can verify quickly with:
\begin{verbatim}
bash scripts/verify_capstone_notebook.sh
\end{verbatim}
